%!TEX root = ED_Grafo.tex

\section{Exemplos}
\begin{frame}
	\frametitle{Exercício}
	\framesubtitle{Algoritmo para Encontrar Amigos em uma Rede Social}
	
	Dada uma rede social representada por um grafo, onde cada aresta representa 
	uma relação de amizade, elabore um algoritmo que liste sugestões de prováveis 
	amigos. \pause Uma pessoa pode ser sugerida a outra se elas não são amigas 
	(na rede social) e têm pelo menos $n$ amigos em comum.
	\pause
	\begin{enumerate}
		\item Suponha o grafo $G(V,A)$ e um vértice $v$ inicial.
		\item Para cada vértice $u$ tal que $^{\exists}(v,u)$:
			\begin{enumerate}
				\item Para cada vértice $w \not\in\{v,u\}$ tal que $^{\exists}(u,w)$:
				\subitem[]{Se $\not\exists (v,w)$: incremente $contador\_W$.}
				\item Se $contador\_W \geq n$: adicione $w$ a lista de sugestões.
			\end{enumerate}
		\item Retorne a lista de sugestões.
	\end{enumerate}
\end{frame}

\begin{frame}[label=notas]
	\frametitle{Exercício}
	\framesubtitle{Encontrar Amigos de $A$}
	
	\begin{center}
		\begin{tikzpicture}[every node/.style={circle,draw},node distance=15mm]
			\node (A) {A};
			\node[above right of=A] (B) {B} edge (A);
			\node[below right of=A] (C) {C} edge (A);
			\node[right of=A] (D) {D} edge (A) edge (B) edge (C);
			\node[right of=D] (E) {E} edge (B) edge (C) edge (D);
			\node[left of=A] (K) {K} edge (A);
			\node[below of=K] (F) {F} edge (A) edge (C);
			\node[above left of=K] (G) {G} edge (B) edge (F);
			\node[above right of=G] (H) {H} edge (B) edge (G);
			\node[above of=B] (I) {I} edge (B) edge (H);
			\node[above of=E] (J) {J} edge (B) edge (E) edge (I);
			\node[above right of=J, right = 2\baselineskip, label=right:{L}] (L) {
\includegraphics[width=.1\textwidth]{foreveralone}};
		\end{tikzpicture}
		\\\vfill		
		\uncover<2->{
		\begin{tabular}{*{13}{|c}|}
			\hline
			& A & B & C & D & E & F & G & H & I & J & K & L\\\hline
			contador & 0 & 0 & 0 & 0 & 3 & 0 & 2 & 1 & 1 & 1 & 0 & 0 \\\hline
		\end{tabular}
		}
	\end{center}
	\begin{columns}[T]
		\column{.45\textwidth}
		\subitem[a)]{Para $n = 3$ \pause $\rightarrow\{E\}$}
		\column{.45\textwidth}
		\subitem[b)]{Para $n = 2$ \pause $\rightarrow\{E, G\}$}
	\end{columns}
\end{frame}

\begin{frame}[label=notas]
	\frametitle{Exemplo: Mapas de Música}
	\framesubtitle{\url{http://audiomap.tuneglue.net/}}
	
	\vfill
	\begin{center}
		\includegraphics[width=.8\textwidth]{musicmap}
		\framezoom<2 | handout:0><2 | handout:0>[border](10mm,30mm)(4cm,4cm)
	\end{center}
\end{frame}